\begin{abstract}
Run-to-failure engine data can make a maintenance warning look successful for
the wrong reason: elapsed cycle count is already informative about how close an
engine is to its recorded failure. This study tests whether measured telemetry
adds useful final-window warning information beyond elapsed cycle and current
operating conditions. We label a simulated turbofan engine cycle positive when
$\mathrm{RUL}\leq20$ and use 90 complete trajectories from the public,
$100\times$-downsampled tiny-N-CMAPSS training partition. The corrected input
definition uses the 14 measured physical channels documented by N-CMAPSS. It
excludes the auxiliary flight-class and health-state fields, unit ID, and RUL.

Five outer engine-held-out folds and four inner engine-grouped folds select a
threshold for a 5\% pre-window false-alert budget. The cycle-plus-operating
baseline recalls 55.4\% of final-window cycles. Adding the 14 measured channels
raises recall to 79.3\% at the matched alert budget. The paired engine-cluster
bootstrap estimates a 23.9-percentage-point recall gain, with a 95\% interval
from 18.3 to 29.6 percentage points.

The result shows that the measured telemetry in this simulator reduces missed
final-window cycles beyond a clock-and-operating-state baseline. It does not
set a service interval, validate a fielded aircraft system, or certify safety.
\end{abstract}

\section{Introduction}

Maintenance teams do not act on a remaining useful life (RUL) estimate alone.
They need a warning policy that identifies a near-term maintenance window while
keeping alerts during earlier operation under control. That makes the relevant
question a classification and threshold-selection problem: does the available
information identify cycles near the decision boundary on engines the model has
not seen?

N-CMAPSS provides simulated aircraft-engine degradation trajectories under
recorded flight conditions~\cite{ariaschao2021ncmapss}. Published work on this
dataset has largely reported RUL regression scores with different readers,
windows, models, and loss functions~\cite{ariaschao2022fusing,lovberg2021variable,solismartin2021stacked}.
Those RMSE or challenge-score results are not directly comparable to this
study, because this study evaluates a binary warning for $\mathrm{RUL}\leq20$
and chooses its decision threshold under a false-alert budget.

Run-to-failure data also require a demanding baseline. An engine's elapsed
cycle number is associated with its remaining life by construction. If all
engines failed at the same cycle, a clock-only rule could solve much of the
task. The trajectories used here fail after different lifetimes, so cycle
number is informative but not sufficient. A useful sensor result must exceed
both the cycle-only and cycle-plus-operating-condition baselines.

This paper first compares four information sets with the same fixed
random-forest reference and the same nested engine-level evaluation: elapsed
cycle only, cycle plus operating conditions, measured telemetry only, and all
observable inputs. It then tests a small candidate set of tabular classifiers
inside the outer-training folds to ask whether the warning policy can improve
without using test-engine labels. The focus remains information value, not an
architecture contest. The claim is narrow: in this simulated setting, the 14
measured channels add final-window warning information beyond elapsed cycle and
recorded operating conditions. The experiment does not establish an aircraft
maintenance schedule.
